Integrable systems occupy a central role in classical and mathematical physics, providing rare examples of nonlinear dynamical systems whose evolution can be described exactly \cite{babelon2003introduction,perelomov1990integrable}. Their origin lies in classical and celestial mechanics, where the search for explicit solutions to the equations of motion led to the identification of systems admitting sufficiently many conserved quantities. This program culminated in the Liouville–Arnold theorem, which gives a geometric characterization of integrability in terms of invariant tori and action–angle variables \cite{babelon2003introduction}.

A decisive conceptual shift occurred in the second half of the twentieth century, when integrability was recognized not merely as solvability by quadratures, but as the manifestation of an underlying algebraic structure. A key development in this direction is the introduction of Lax pairs, which encode the dynamics in a matrix equation of the form $\dot{L} = [M,L]$ \cite{https://doi.org/10.1002/cpa.3160210503}. This formulation, originally introduced in the study of nonlinear evolution equations such as the Korteweg–de Vries equation \cite{gardner1967method,Zakharov1971}, allows one to interpret the time evolution as an isospectral deformation and to construct conserved quantities from the spectral invariants of the Lax matrix.

This algebraic viewpoint is further refined by the $r$-matrix formalism, which organizes the Poisson brackets of the entries of the Lax matrix into a tensorial structure \cite{SemenovTyanShanskii1983,Sklyanin1982}. In this framework, the involutivity of the spectral invariants follows from a precise algebraic relation involving the $r$-matrix. These structures reveal a deep connection between integrable systems and Lie algebras, with the tensor Casimir element and related invariant tensors playing a central role \cite{perelomov1990integrable}.

The aim of this work is to present these ideas in a concrete and largely self-contained manner, with particular emphasis on the interplay between Hamiltonian dynamics, Lax representations, and algebraic structures. After reviewing Liouville integrability and the basic properties of Lax pairs, we develop the $r$-matrix formalism, discuss the role of the Casimir element, and derive the classical and modified Yang–Baxter equations. These concepts are then illustrated on the example of the open Toda chain \cite{toda1967vibration,Olshanetsky1979,KOSTANT1979195}.

While Sections~\ref{subsec:Introduction}--\ref{subsec:Action--angle variables} and Section~\ref{sec:Commuting flows} follow standard presentations (mainly based on \cite{babelon2003introduction}), several parts of this work are developed independently. In particular, Section~\ref{subsec:Example: Lax pairs for oscillators} provides a self-contained derivation of a Lax representation for one-dimensional systems, while Sections~\ref{sec:rmatrix_formalism}, \ref{sec:casimir}, and \ref{sec:YB} present a direct and explicit treatment of the $r$-matrix formalism, the tensor Casimir element, and the classical and modified Yang–Baxter equations.

The proof of the $r$-matrix structure for the open Toda chain in Section~\ref{sec:toda_rmatrix} is given in detail, and Section~\ref{sec:Independence of Hk} contains a direct argument for the independence of the Hamiltonians. The goal is to provide a coherent and partially self-contained presentation, with emphasis on explicit constructions and derivations.